% ============================================================
%  Part V — Verbs
% ============================================================
\gpart{V}{Verbs}

% ------------------------------------------------------------
\gsec{Present tense — regular endings}

\gintro{German has a single present tense covering all three English forms, so \textit{ich spiele} can mean \emph{I play}, \emph{I am playing} or \emph{I do play}.}

Take the infinitive, remove \textit{-en}, add the ending.
\textit{spielen} $\rightarrow$ stem \textit{spiel-}

\begin{gtbl}{colspec={X[0.85,l] X[0.75,c] X[1,l] X[0.7,c] X[1.1,l]}}
  Number & Person & Pronoun & Ending & spielen \\
  Singular & 1st & ich & \en{-e}  & spiele  \\
  Singular & 2nd & du  & \en{-st} & spielst \\
  Singular & 3rd & er / sie / es & \en{-t} & spielt \\
  Plural   & 1st & wir & \en{-en} & spielen \\
  Plural   & 2nd & ihr & \en{-t}  & spielt  \\
  Plural   & 3rd & sie / Sie & \en{-en} & spielen \\
\end{gtbl}

\begin{gnote}
There are only four distinct endings and three of them repeat, so the whole
present tense of a regular verb is really \textit{-e, -st, -t, -en}. The formal
\textit{Sie} always uses the 3rd person plural form.
\end{gnote}

\tcap{Spelling adjustments}

\begin{gtbl}{colspec={X[1,l] X[0.9,l] X[1.1,l]}}
  Stem ends in & Change & Example \\
  -d, -t, -chn, -ffn, -gn & insert \en{e} before -st and -t & du arbeit\textbf{e}st, er arbeit\textbf{e}t \\
  -s, -ß, -z, -x, -tz & du takes only \en{-t} & du hei\textbf{ßt}, du tan\textbf{zt} \\
  -eln, -ern & 1st sg may drop the e & ich samm\textbf{le}, ich änd\textbf{re} \\
\end{gtbl}

\tcap{German has one present tense}

\begin{flattbl}{colspec={X[1,l] X[1.7,l]}}
  ich spiele & I play \textbullet{} I am playing \textbullet{} I do play \\
  Add \textit{gerade} & to stress that it is happening right now \\
\end{flattbl}

% ------------------------------------------------------------
\gsec{Stem-changing verbs}

\gintro{Around sixty common verbs shift their stem vowel, but only in two of the six present-tense forms. The plural is always regular.}

Many strong verbs change their stem vowel in the \textbf{du} and \textbf{er/sie/es}
forms only. Everything else stays regular.

\begin{gtbl}{colspec={X[0.7,l] X[0.95,l] X[0.85,l] X[0.85,l] X[1,l]}}
  Change & Infinitiv & du & er/sie/es & Also \\
  a $\rightarrow$ ä  & fahren  & fährst  & fährt  & schlafen, tragen, halten \\
  a $\rightarrow$ ä  & fallen  & fällst  & fällt  & lassen, waschen, gefallen \\
  au $\rightarrow$ äu & laufen & läufst  & läuft  & saufen \\
  e $\rightarrow$ i  & geben   & gibst   & gibt   & helfen, sprechen, essen \\
  e $\rightarrow$ i  & treffen & triffst & trifft & werfen, brechen, sterben \\
  e $\rightarrow$ ie & sehen   & siehst  & sieht  & lesen, empfehlen, stehlen \\
  irregular & nehmen & nimmst & nimmt & \\
  irregular & werden & wirst  & wird  & \\
  irregular & wissen & weißt  & weiß  & \\
\end{gtbl}

\begin{gnote}
The stem change never appears in the plural: \textit{wir fahren},
\textit{ihr fahrt}, \textit{sie fahren}. It also disappears in the
\textit{du} imperative of a\,$\rightarrow$\,ä verbs (\textit{Fahr!}) but is kept
for e\,$\rightarrow$\,i verbs (\textit{Gib!}, \textit{Sprich!}).
\end{gnote}

% ------------------------------------------------------------
\gsec{haben, sein, werden}

\gintro{These three carry every compound tense in the language: \textit{haben} and \textit{sein} build the perfect tenses, \textit{werden} builds the future and the passive.}

\tcap{Präsens}

\begin{gtbl}{colspec={X[1.25,l] X[1,c] X[1,c] X[1,c]}}
  Person & haben & sein & werden \\
  ich       & habe  & bin  & werde  \\
  du        & hast  & bist & wirst  \\
  er/sie/es & hat   & ist  & wird   \\
  wir       & haben & sind & werden \\
  ihr       & habt  & seid & werdet \\
  sie/Sie   & haben & sind & werden \\
\end{gtbl}

\tcap{Präteritum \eng{simple past}}

\begin{gtbl}{colspec={X[1.25,l] X[1,c] X[1,c] X[1,c]}}
  Person & haben & sein & werden \\
  ich       & hatte   & war   & wurde   \\
  du        & hattest & warst & wurdest \\
  er/sie/es & hatte   & war   & wurde   \\
  wir       & hatten  & waren & wurden  \\
  ihr       & hattet  & wart  & wurdet  \\
  sie/Sie   & hatten  & waren & wurden  \\
\end{gtbl}

\tcap{The rest of the paradigm}

\begin{gtbl}{colspec={X[1.25,l] X[1,c] X[1,c] X[1,c]}}
  Form & haben & sein & werden \\
  Partizip II   & gehabt & gewesen & geworden \\
  Konjunktiv II & hätte  & wäre    & würde \\
  Perfekt with  & haben  & sein    & sein \\
\end{gtbl}

\begin{gnote}
\textit{werden} does three separate jobs: \emph{to become}
(\textit{Er wird Arzt}), the future auxiliary (\textit{Ich werde gehen}) and the
passive auxiliary (\textit{Das Haus wird gebaut}). In the passive perfect its
participle shortens to \textit{worden}.
\end{gnote}

% ------------------------------------------------------------
\gsec{Modal verbs}

\gintro{Modal verbs express permission, ability, obligation and intention. They take a second verb in the infinitive, which moves to the end of the clause.}

\begin{gtbl}{colspec={X[1,l] X[1,c] X[1,c] X[1,c] X[1,c] X[1,c] X[1,c]}, rows={font=\footnotesize}}
  Person & dürfen & können & mögen & müssen & sollen & wollen \\
  ich       & darf   & kann   & mag   & muss  & soll   & will   \\
  du        & darfst & kannst & magst & musst & sollst & willst \\
  er/sie/es & darf   & kann   & mag   & muss  & soll   & will   \\
  wir       & dürfen & können & mögen & müssen & sollen & wollen \\
  ihr       & dürft  & könnt  & mögt  & müsst & sollt  & wollt  \\
  sie/Sie   & dürfen & können & mögen & müssen & sollen & wollen \\
\end{gtbl}

\begin{gnote}
Two things are true of every modal. The \textbf{ich and er forms are identical
and take no ending}, and the singular loses the umlaut of the infinitive.
\textit{sollen} and \textit{wollen} never had an umlaut to lose; of those two
only \textit{sollen} keeps its vowel unchanged, while \textit{wollen} shifts
o\,$\rightarrow$\,i.
\end{gnote}

\tcap{Meaning}

\begin{flattbl}{colspec={X[0.8,l] X[1,l] X[1.6,l]}}
  dürfen & may, be allowed to & Darf ich rauchen? \\
  können & can, be able to    & Ich kann schwimmen. \\
  mögen  & to like            & Ich mag Kaffee. \\
  müssen & must, have to      & Ich muss gehen. \\
  sollen & should, be supposed to & Du sollst warten. \\
  wollen & to want to         & Er will kommen. \\
  möchten & would like to \eng{polite} & Ich möchte einen Tee. \\
\end{flattbl}

\begin{gnote}
The modal is conjugated in second position and the \textbf{main verb goes to the
end as a bare infinitive}: \textit{Ich \textbf{muss} morgen früh
\textbf{aufstehen}}. Negated \textit{müssen} means \emph{need not}, not
\emph{must not} — for that use \textit{nicht dürfen}.
\end{gnote}

% ------------------------------------------------------------
\gsec{Separable and inseparable verbs}

\gintro{A prefix changes a verb's meaning. Whether it detaches from the stem depends on where the stress falls, which is why the two groups behave so differently.}

\tcap{Separable — prefix is stressed and detaches}

\begin{flattbl}{colspec={X[1,l] X[1,l] X[1,l] X[1,l]}}
  ab-   & an-   & auf-  & aus- \\
  bei-  & ein-  & mit-  & nach- \\
  vor-  & zu-   & zurück- & weg- \\
  hin-  & her-  & los-  & fest- \\
\end{flattbl}

\begin{gtbl}{colspec={X[0.9,l] X[2.1,l]}}
  Position & aufstehen \eng{to get up} \\
  Main clause    & Ich \textbf{stehe} um sechs \textbf{auf}. \\
  Subordinate    & \dots, weil ich um sechs \textbf{aufstehe}. \\
  Partizip II    & Ich bin um sechs \textbf{aufgestanden}. \\
  zu + infinitive & Ich versuche, früh \textbf{aufzustehen}. \\
\end{gtbl}

\tcap{Inseparable — prefix is unstressed and never moves}

\begin{flattbl}{colspec={X[1,l] X[1,l] X[1,l] X[1,l]}}
  be-  & emp-  & ent-  & er- \\
  ge-  & miss- & ver-  & zer- \\
\end{flattbl}

\begin{gnote}
Inseparable verbs form the \textbf{Partizip II without ge-}:
\textit{verstehen} $\rightarrow$ \textit{verstanden},
\textit{bezahlen} $\rightarrow$ \textit{bezahlt}. The same applies to verbs
ending in \textit{-ieren}: \textit{studieren} $\rightarrow$ \textit{studiert}.
\end{gnote}

\tcap{Dual prefixes — stress decides}

\begin{flattbl}{colspec={X[0.6,l] X[1.2,l] X[1.2,l]}}
  über- & \textbf{über}setzen \eng{ferry across, sep.} & über\textbf{setz}en \eng{translate, insep.} \\
  unter- & \textbf{unter}gehen \eng{sink, sep.} & unter\textbf{schrei}ben \eng{sign, insep.} \\
  um-   & \textbf{um}ziehen \eng{move house, sep.} & um\textbf{ar}men \eng{embrace, insep.} \\
  durch- & \textbf{durch}fallen \eng{fail, sep.} & durch\textbf{schau}en \eng{see through, insep.} \\
\end{flattbl}
